%!TEX root = forallxyyc.tex
\part{一阶逻辑的自然演绎}
\label{ch.NDFOL}
\addtocontents{toc}{\protect\mbox{}\protect\hrulefill\par}

\chapter{一阶逻辑的基本规则}\label{s:BasicFOL}

一阶逻辑的语言使用了真值函项逻辑的所有联结词。因此，一阶逻辑中的证明也将使用\cref{ch.NDTFL}中的所有基本规则和派生规则。我们也将沿用那里引入的证明论概念（特别是符号$\proves$）。不过，我们还需要一些新的基本规则来支配量词以及等号。

\section{全称消去}\label{s:forallElim}

从“所有东西都是~$F$”这个主张，你可以推出任何特定的东西是~$F$。你说出它的名称；它就是~$F$。所以下面这个推导应该是没问题的：


\begin{fitchproof}
	\hypo{a}{\forall x\,\atom{R}{x,x,d}}\PR
	\have{c}{\atom{R}{a,a,d}} \Ae{a}
\end{fitchproof}


我们通过去掉全称量词并把所有$x$的出现都替换为$a$得到了第~$2$行。同样地，下面的推导也应该被允许：


\begin{fitchproof}
	\hypo{a}{\forall x\,\atom{R}{x,x,d}}\PR
	\have{c}{\atom{R}{d,d,d}} \Ae{a}
\end{fitchproof}


我们这里通过去掉全称量词并把所有$x$的出现都替换为$d$得到了第~$2$行。我们也可以用任何其他我们想用的名称来做同样的事。

这就引出了全称消去规则（$\forall$E）：

\factoidbox{


\begin{fitchproof}
	\have[m]{a}{\forall \metav{x}\,\metav{A}(\ldots \metav{x} \ldots \metav{x}\ldots)}
	\have[\ ]{c}{\metav{A}(\ldots \metav{c} \ldots \metav{c}\ldots)} \Ae{a}
\end{fitchproof}

}

这里的记法在\cref{s:TruthFOL}中介绍过。关键在于，你可以得到全称量化公式的任何\emph{代入实例}：把被量化变元的每一个自由出现都替换为你喜欢的任何名称。

我们应该强调，（与所有消去规则一样）只有当全称量词是主逻辑算子时，你才能应用 $\forall$E 规则。因此下面的推导是\emph{被禁止的}：


\begin{fitchproof}
	\hypo{a}{\forall x\,\atom{B}{x} \eif \atom{B}{k}}\PR
	\have{c}{\atom{B}{b} \eif \atom{B}{k}}\by{违规调用 $\forall$E 的尝试}{a}
\end{fitchproof}


这是不合法的，因为$\forall x$不是第~$1$行的主逻辑算子。（如果你需要提醒为何这类推理应该被禁止，请重读\cref{s:MoreMonadic}。）

\section{存在引入}\label{s:existsIntro}

从“某个特定的东西是~$F$”这个主张，你可以推出有东西是~$F$。所以我们应该允许：


\begin{fitchproof}
	\hypo{a}{\atom{R}{a,a,d}}\PR
	\have{b}{\exists x\, \atom{R}{a,a,x}} \Ei{a}
\end{fitchproof}


这里，我们把名称$d$替换为变元$x$，然后对其存在量化。同样地，我们也会允许：


\begin{fitchproof}
	\hypo{a}{\atom{R}{a,a,d}}\PR
	\have{c}{\exists x\, \atom{R}{x,x,d}} \Ei{a}
\end{fitchproof}


这里我们把名称$a$的两次出现都替换为变元，然后进行了存在概括。但我们并不需要把名称的\emph{两次}出现都替换为变元：如果那喀索斯爱自己，那么就存在某个爱那喀索斯的人。所以我们也允许：


\begin{fitchproof}
	\hypo{a}{\atom{R}{a,a,d}}\PR
	\have{d}{\exists x\, \atom{R}{x,a,d}} \Ei{a}
\end{fitchproof}


这里我们先把名称$a$的\emph{一个}出现替换为变元，然后进行了存在概括。这些观察引出了我们的引入规则，不过为了解释它，我们需要引入一些新的记法。

当句子 $\metav{A}$ 包含名称 $\metav{c}$ 时，我们可以通过写作$\metav{A}(\ldots \metav{c} \ldots
\metav{c}\ldots)$来强调这一点。我们将用$\metav{A}(\ldots \metav{x} \ldots
\metav{c}\ldots)$来表示通过将名称 \metav{c} 的\emph{部分或全部}出现替换为变元~\metav{x} 而得到的任何公式（前提是，无论在何处替换，$\metav{c}$ 都不出现在约束~$\metav{x}$ 的量词的辖域内）。有了这些准备，我们的引入规则 $\exists$I 是：
\factoidbox{


\begin{fitchproof}
	\have[m]{a}{\metav{A}(\ldots \metav{c} \ldots \metav{c}\ldots)}
	\have[\ ]{c}{\exists \metav{x}\,\metav{A}(\ldots \metav{x} \ldots \metav{c}\ldots)} \Ei{a}
\end{fitchproof}


%\metav{x} must not occur in $\metav{A}(\ldots \metav{c} \ldots
%\metav{c}\ldots)$
}
% The constraint is included to guarantee that any application of the rule yields a  sentence of FOL. Thus the following is allowed:
% \begin{fitchproof}
% 	\hypo{a}{\atom{R}{a,a,d}}
% 	\have{d}{\exists x\, \atom{R}{x,a,d}} \Ei{a}
% 	\have{e}{\exists y \exists x\, \atom{R}{x,y,d}} \Ei{d}
% \end{fitchproof}
% But this is banned:
% \begin{fitchproof}
% 	\hypo{a}{\atom{R}{a,a,d}}
% 	\have{d}{\exists x\, \atom{R}{x,a,d}} \Ei{a}
% 	\have{e}{\exists x\, \exists x\, \atom{R}{x,x,d}}\by{naughty attempt to invoke $\exists$I}{d}
% \end{fitchproof}
% since the expression on line~3 contains clashing variables, and so is not a sentence of FOL.

\section{空论域}
下面的证明结合了我们关于量词的两条新规则：


\begin{fitchproof}
		\hypo{a}{\forall x\, \atom{F}{x}}\PR
		\have{in}{\atom{F}{a}}\Ae{a}
		\have{e}{\exists x\, \atom{F}{x}}\Ei{in}
	\end{fitchproof}


这会是一个糟糕的证明吗？如果确实存在任何东西，那么我们当然可以从“所有东西都是$F$”这个事实，推断出“有东西是$F$”。但如果\emph{根本没有}东西存在呢？那么，“所有东西都是$F$”显然就是\emph{空真}的；然而，这并不能推出“有东西是$F$”，因为没有什么东西可以去\emph{成为}$F$。所以，如果我们声称，单凭逻辑，$\exists x\,\atom{F}{x}$就从$\forall x\,\atom{F}{x}$得出，那么我们就是在声称，单凭\emph{逻辑}，就存在某物而非无。这可能会让我们觉得有点奇怪。

事实上，我们已经接受了这种怪异之处。在\cref{s:FOLBuildingBlocks}中，我们规定，一阶逻辑中的论域必须至少有一个成员。然后，我们把（一阶逻辑的）有效性定义为一个在所有解释下都为真的语句。由于$\exists x\, x=x$将在每个解释下都为真，这\emph{同样}起到了规定作用：逻辑本身决定了存在某物而非无。

既然逻辑是否应该告诉我们“必须有某物而非无”这一点远非显而易见，我们在这里可能确实有点作弊的嫌疑。

然而，如果我们拒绝作弊，代价就很高。下面是我们想坚持的三点：


\begin{compactlist}
		\item $\forall x\,\atom{F}{x} \proves \atom{F}{a}$：毕竟，那就是 $\forall$E 规则。
		\item $\atom{F}{a} \proves \exists x\,\atom{F}{x}$：毕竟，那就是 $\exists$I 规则。
		\item 将证明拼接起来的能力：毕竟，推理就是通过将许多小步骤组合成相当长的链条来运作的。
	\end{compactlist}


如果我们在上述三个方面都如愿以偿，那么我们就必须承认$\forall x\,\atom{F}{x} \proves \exists x\,\atom{F}{x}$。因此，如果我们想要这三个特性，证明系统本身就告诉我们存在某物而非无。而如果我们拒绝接受这一点，那么我们就得放弃我们想要坚持的三个特性之一！

在我们开始考虑放弃哪一个之前，我们不妨问问这个“作弊”究竟有多严重。诚然，它可能会使我们更难参与关于“为何存在某物而非无”的神学辩论。但在其余情况下，一切都会相安无事。所以，也许我们只需把我们的证明系统（以及更广义的一阶逻辑）看作具有一个略微受限的适用范围。如果我们将来想要允许“无”的可能性，那我们就不得不寻找一个更复杂的证明系统。但只要我们满足于忽略这种可能性，我们的证明系统就完全正常。（同样，每个论域必须包含至少一个对象的规定也是如此。）

\section{全称引入}\label{s:forallIntro}

假设你已经证明了每一个特定的东西都是 F（并且没有其他东西需要考虑）。那么，你就有理由声称所有东西都是 F。这会启发我们得到以下证明规则。如果你已经确立了$\forall x\,\atom{F}{x}$的每一个代入实例，那么你就可以推断出$\forall x\,\atom{F}{x}$。

不幸的是，这条规则将完全无法使用。要确立每一个单独的代入实例，就需要证明$\atom{F}{a}$、$\atom{F}{b}$、……、$\atom{F}{j_2}$、……、$\atom{F}{r_{79002}}$、……，如此等等。事实上，由于一阶逻辑中有无穷多个名称，这个过程永无止境。所以我们永远无法应用那条规则。我们需要更巧妙一点，来制定我们的全称引入规则。

考虑以下论证会带来启发：
$$\forall x\,\atom{F}{x} \therefore \forall y\,\atom{F}{y}$$
这个论证应该是\emph{显然}有效的。毕竟，字母的变化应该只是个人偏好问题，而没有任何逻辑后果。但是我们的证明系统如何体现这一点呢？假设我们这样开始一个证明：


\begin{fitchproof}
	\hypo{x}{\forall x\, \atom{F}{x}} \PR
	\have{a}{\atom{F}{a}} \Ae{x}
\end{fitchproof}


我们已经证明了$\atom{F}{a}$。当然，没有什么能阻止我们用同样的证明依据去证明$\atom{F}{b}$、$\atom{F}{c}$、……、$\atom{F}{j_2}$、……、$\atom{F}{r_{79002}}$……直到我们耗尽空间、时间或耐心。但仔细想想，我们发现，对于任何名称~\metav{c}，我们都有办法证明 $\atom{F}{\metav{c}}$。如果我们能对\emph{任何}东西做到这一点，我们当然应该说$F$对\emph{所有}东西都成立。因此，这为我们如下推出$\forall y\,\atom{F}{y}$提供了正当理由：


\begin{fitchproof}
	\hypo{x}{\forall x\, \atom{F}{x}}\PR
	\have{a}{\atom{F}{a}} \Ae{x}
	\have{y}{\forall y\, \atom{F}{y}} \Ai{a}
\end{fitchproof}


这里的关键想法是，$a$只是一个\emph{任意}的名称。它没有什么特别的——我们本可以选择任何其他名称——而且证明依然成立。而这个关键想法启发了全称引入规则（$\forall$I）：
\factoidbox{

\begin{fitchproof}
	\have[m]{a}{\metav{A}(\ldots \metav{c} \ldots \metav{c}\ldots)}
	\have[\ ]{c}{\forall \metav{x}\,\metav{A}(\ldots \metav{x} \ldots \metav{x}\ldots)} \Ai{a}
\end{fitchproof}

其中，名称 \metav{c} 不得出现在前提或未解除的假设中。%\\
%\metav{x} must not occur in $\metav{A}(\ldots \metav{c} \ldots
%\metav{c}\ldots)$
}

在\cref{s:MainLogicalOperatorQuantifier}中，我们引入了一个用于指示\emph{代入实例}的约定：我们说，如果$\metav{A}(\ldots \metav{x} \ldots \metav{x}\ldots)$是一个公式，那么$\metav{A}(\ldots \metav{c} \ldots \metav{c}\ldots)$就是将 $\metav{A}$ 中\emph{每一个}$\metav{x}$ 的自由出现替换为~$\metav{c}$ 的结果。为了陈述我们的 $\forall$I 规则，我们使用了同一约定，只是方向相反：如果$\metav{A}(\ldots \metav{c} \ldots \metav{c}\ldots)$是一个公式，那么$\metav{A}(\ldots \metav{x} \ldots \metav{x}\ldots)$就是将其中\emph{每一个}~$\metav{c}$ 的出现替换为~$\metav{x}$ 的结果。（仅当 $\metav{c}$ 在 $\metav{A}$ 中未出现在约束~$\metav{x}$ 的量词的辖域内时，才允许这样做。）

在应用 $\forall$I 规则时，有两点需要特别注意。如果我们在选择名称~$\metav{c}$ 并将其替换为~$\metav{x}$ 时不小心，就有可能做出不正确的推理。

$\forall$I 规则的限制条件要求 \metav{c} 不得出现在前提或未解除的假设中。这个限制确保了我们始终在足够一般的层面上进行推理。%\footnote{Recall from\cref{s:BasicTFL}that we are treating `$\ered$' as a canonical contradiction. But if it were the canonical contradiction as involving some \emph{constant}, it might interfere with the constraint mentioned here. To avoid such problems, we will treat `$\ered$' as a canonical contradiction \emph{that involves no particular names}.}
为了看看这个限制是如何起作用的，考虑下面这个糟糕的论证：

\begin{earg}
		\item 人人都爱 Kylie Minogue。
		\item[\texttherefore] 人人都爱自己。
	\end{earg}

我们可能将这个明显无效的推理模式符号化为：
$$\forall x\,\atom{L}{x,k} \therefore \forall x\,\atom{L}{x,x}$$
假设我们试图提供一个证明来支持这一论证：

\begin{fitchproof}
	\hypo{x}{\forall x\, \atom{L}{x,k}}\PR
	\have{a}{\atom{L}{k,k}} \Ae{x}
	\have{y}{\forall x\, \atom{L}{x,x}} \by{试图违规调用 $\forall$I}{a}
\end{fitchproof}

这是不允许的，因为$k$已经出现在一个前提中，即第 $1$ 行。关键在于，如果我们对正在处理的对象做出了任何假设，那么我们的推理就不够一般，从而无法允许使用 $\forall$I。

不过，虽然名称不得出现在任何\emph{未解除的}假设中，但它可以出现在\emph{已解除的}假设中。也就是说，它可以出现在我们已经关闭的子证明里。例如，下面这样是完全没问题的：

\begin{fitchproof}
	\open
		\hypo{f1}{\atom{G}{d}}\AS
		\have{f2}{\atom{G}{d}}\by{R}{f1}
	\close
	\have{ff}{\atom{G}{d} \eif \atom{G}{d}}\ci{f1-f2}
	\have{zz}{\forall z(\atom{G}{z} \eif \atom{G}{z})}\Ai{ff}
\end{fitchproof}

这告诉我们，$\forall z (\atom{G}{z} \eif \atom{G}{z})$是一个\emph{定理}。理应如此。

如果我们只替换\emph{某些}名称而不替换其他的，就会“证明”出一些荒谬的事情。例如，考虑这个论证：

\begin{earg}
	\item 每个人都和自己一样老。
	\item[\texttherefore] 每个人都和 Judi Dench 一样老。
	\end{earg}

我们可能会这样符号化：
$$\forall x\,\atom{O}{x,x} \therefore \forall x\,\atom{O}{x,d}$$
但现在假设我们试图用下面的推导来支持这个糟糕的论证：

\begin{fitchproof}
	\hypo{x}{\forall x\, \atom{O}{x,x}}\PR
	\have{a}{\atom{O}{d,d}}\Ae{x}
	\have{y}{\forall x\, \atom{O}{x,d}}\by{试图违规调用 $\forall$I}{a}
\end{fitchproof}

幸运的是，我们的规则不允许我们这样做：这个尝试的证明被禁止了，因为它并没有将第 $2$ 行中\emph{每一个}$d$的出现都替换为$x$。

\section{存在量词消去}\label{s:existsElim}

假设我们知道\emph{某物}是~$F$。问题在于，仅仅知道这一点并不能告诉我们哪个东西是~$F$。所以，看起来从$\exists x\,\atom{F}{x}$出发，我们不能直接得出$\atom{F}{a}$、$\atom{F}{e_{23}}$或该语句的任何其他代入实例。我们能做什么呢？

假设我们知道有某物是~$F$，并且所有是~$F$ 的东西也都是~$G$。用（近似）自然的英语，我们可能会这样推理：

\begin{quote}
		既然有某物是 $F$，那就存在某个特定的东西是 $F$。除了知道它是一个 $F$ 之外，我们对它一无所知，但为了方便起见，不妨称它为“贝姬”。所以：贝姬是 $F$。因为所有是 $F$ 的东西都是~$G$，所以贝姬是~$G$。但既然贝姬是~$G$，那就说明有某物是~$G$。而且，这个结论并不依赖于贝姬具体是哪个对象。因此，有某物是~$G$。
	\end{quote}

我们可能会尝试在证明中捕捉这个推理模式，如下所示：

\begin{fitchproof}
	\hypo{es}{\exists x\, \atom{F}{x}}\PR
	\hypo{ast}{\forall x(\atom{F}{x} \eif \atom{G}{x})}\PR
	\open
		\hypo{s}{\atom{F}{b}}\AS
		\have{st}{\atom{F}{b} \eif \atom{G}{b}}\Ae{ast}
		\have{t}{\atom{G}{b}} \ce{st, s}
		\have{et1}{\exists x\, \atom{G}{x}}\Ei{t}
	\close
	\have{et2}{\exists x\, \atom{G}{x}}\Ee{es,s-et1}
\end{fitchproof}

\noindent
拆解一下这个过程：我们先写下了前提。在第 $3$ 行，我们增加了一个额外的假设：$\atom{F}{b}$。这只是$\exists x\,\atom{F}{x}$的一个代入实例。基于这个假设，我们得出了$\exists x\,\atom{G}{x}$。注意，对于名称$b$所指的对象，我们并没有做任何\emph{特殊}假设；我们\emph{只}假设了它满足$\atom{F}{x}$。因此，推理过程不依赖于它具体是哪个对象。而第 $1$ 行告诉我们\emph{有东西}满足$\atom{F}{x}$，所以我们的推理模式是完全一般的。我们可以解除这个具体的假设$\atom{F}{b}$，只单独推出$\exists x\,\atom{G}{x}$。

综上所述，我们得到了存在量词消去规则（$\exists$E）：
\factoidbox{


\begin{fitchproof}
	\have[m]{a}{\exists \metav{x}\,\metav{A}(\ldots \metav{x} \ldots \metav{x}\ldots)}
	\open
		\hypo[i]{b}{\metav{A}(\ldots \metav{c} \ldots \metav{c}\ldots)}\AS
		\have[j]{c}{\metav{B}}
	\close
	\have[\ ]{d}{\metav{B}} \Ee{a,b-c}
\end{fitchproof}


\metav{c} 不得出现在第 $i$ 行之前任何未解除的假设中\\
\metav{c} 不得出现在 $\exists \metav{x}\,\metav{A}(\ldots \metav{x} \ldots \metav{x}\ldots)$ 中\\
\metav{c} 不得出现在 \metav{B} 中}
与全称引入规则一样，这些约束条件极其重要。要知道为什么，来看看下面这个糟糕的论证：

\begin{earg}
		\item Tim Button 是讲师。
		\item 有人不是讲师。
		\item[\texttherefore] Tim Button 既是讲师又不是讲师。
	\end{earg}


我们可以将这个显然无效的推理模式符号化如下：
$$\atom{L}{b}, \exists x\, \enot\atom{L}{x} \therefore \atom{L}{b} \eand \enot \atom{L}{b}$$
现在，假设我们试图给出一个证明来支持此论证：

\begin{fitchproof}
	\hypo{f}{\atom{L}{b}}\PR
	\hypo{nf}{\exists x\, \enot \atom{L}{x}}\PR
	\open
		\hypo{na}{\enot \atom{L}{b}}\AS
		\have{con}{\atom{L}{b} \eand \enot \atom{L}{b}}\ai{f, na}
	\close
\ifHTMLtarget
	\have{econ1}{\atom{L}{b} \eand \enot \atom{L}{b}}\by{不当尝试调用 $\exists$E }{nf, na-con}
\else
	\have{econ1}{\atom{L}{b} \eand \enot \atom{L}{b}}\by{不当尝试}{}
	\have[\ ]{x}{}\by{调用 $\exists$E }{nf, na-con}
\fi
\end{fitchproof}


这个证明的最后一行是不允许的。我们在第 $3$ 行用作$\exists x\, \enot \atom{L}{x}$的代入实例中的名称，即$b$，出现在了第 $4$ 行中。下面这样也好不到哪去：

\begin{fitchproof}
	\hypo{f}{\atom{L}{b}}\PR
	\hypo{nf}{\exists x\, \enot \atom{L}{x}}\PR
	\open
		\hypo{na}{\enot \atom{L}{b}}\AS
		\have{con}{\atom{L}{b} \eand \enot \atom{L}{b}}\ai{f, na}
		\have{con1}{\exists x (\atom{L}{x} \eand \enot \atom{L}{x})}\Ei{con}
	\close
\ifHTMLtarget
	\have{econ1}{\exists x (\atom{L}{x} \eand \enot
	\atom{L}{x})}\by{不当尝试调用 $\exists$E }{nf, na-con1}
\else
	\have{econ1}{\exists x (\atom{L}{x} \eand \enot \atom{L}{x})}\by{不当尝试}{}
	\have[\ ]{x}{}\by{调用 $\exists$E }{nf, na-con1}
\fi
\end{fitchproof}


最后一行仍然是不允许的。因为用于$\exists x\, \enot \atom{L}{x}$的代入实例中的名称，即$b$，出现在了一个未解除的假设中，即第 $1$ 行。

最后，看看这个“证明”：

\begin{fitchproof}
	\hypo{nf}{\exists x\, \atom{L}{a, x}}\PR
	\open
		\hypo{na}{\atom{L}{a, a}}\AS
		\have{con1}{\exists x\,\atom{L}{x,x}}\Ei{na}
	\close
\ifHTMLtarget
	\have{econ1}{\exists x\,\atom{L}{x,x}}\by{不当尝试调用 $\exists$E }{nf, na-con1}
\else
	\have{econ1}{\exists x\,\atom{L}{x,x}}\by{不当尝试}{}
	\have[\ ]{x}{}\by{调用 $\exists$E }{nf, na-con1}
\fi
\end{fitchproof}


这里的名称$a$违反了 $\exists$E 规则的第二个条件：它出现在了$\exists x\, \atom{L}{a, x}$中，即第 $1$ 行。我们必须阻止它成为证明，因为它是无效的：“有人喜欢自己”并不能从“Alex 喜欢某人”推出。

这个故事给我们的教训是：\emph{如果你想从存在量词中提取信息，请为你的代入实例选择一个新的名称。}这样你就能保证满足 $\exists$E 规则的所有约束条件。

\section{量词规则与空量化}

我们的量词规则中有一个微妙之处，这与我们语法规则的一个方面有关，即\emph{空量化}的可能性（见\cref{s:FreeVars}末尾）。在\cref{s:forallElim,s:existsElim}中，我们使用了\cref{s:TruthFOL}的替换记号：$\metav{A}(\ldots
\metav{c} \ldots \metav{c} \ldots)$ 是将 $\metav{A}$ 中\emph{每一个}$\metav{x}$ 的\emph{自由}出现替换为~$\metav{c}$ 后得到的公式。这里，我们强调“自由”是有原因的。这意味着你不能总是替换\emph{每一个}~$\metav{x}$ 的出现。例如，下面这样做就是错误的：

\begin{fitchproof}
	\hypo{a}{\forall x(\atom{F}{x} \land \exists x\,\atom{G}{x})}\PR
	\ifHTMLtarget
	\have{e}{\atom{F}{a} \land \exists x\,\atom{G}{a}}\by{不当的尝试：意图应用 $\forall$E}{a}
\else
	\have{e}{\atom{F}{a} \land \exists x\,\atom{G}{a}}\by{不当的尝试}{}
	\have[\ ]{x}{}\by{意图应用 $\forall$E}{a}
\fi
\end{fitchproof}


在第 $2$ 行，我们错误地构造了代入实例：我们把 $\atom{F}{x}$ 和 $\atom{G}{x}$ 中的 $x$ 都替换成了 $a$，但只有前一个 $x$ 是自由出现——后一个被 $\exists x$ 约束了。这不仅会违反替换规则的限制，还会让我们能够“证明”一个无效的论证。由于 $\exists x\,\atom{G}{a}$ 是空量化的情况，它等价于 $\atom{G}{a}$。此外，第 $1$ 行等价于 $\forall x(\atom{F}{x} \land \exists y\,\atom{G}{y})$。
不难看出，这并不蕴涵 $\atom{F}{a} \land
\atom{G}{a}$。你不太可能遇到这样的情况（除非你的老师很刁钻）：你的练习题很可能只涉及像 $\forall x(\atom{F}{x} \land \exists
y\,\atom{G}{y})$ 这样的语句，其中的变元都被清晰地分开了。

在\cref{s:existsIntro}中，我们说过，我们用 $\metav{A}(\ldots
\metav{x} \ldots \metav{c}\ldots)$ 来表示任何通过如下方式获得的公式：将名称 \metav{c} 的部分或全部出现替换为变元~\metav{x}，\emph{但条件是：$\metav{c}$ 在每次被替换时，都不出现在约束 $\metav{x}$ 的量词的辖域内}。这一强调的限制在某些情况下禁止替换 $\metav{c}$，这一点有些微妙。在实践中，你可能永远不会发现自己处于可能违反它的境地。但为了让规则正确，这个限制是必要的。例如，由于我们允许空量化，$\exists
x\forall x\,\atom{L}{x,x}$ 是一阶逻辑的合法语句。只不过，其中的第一个 $\exists x$ 不起作用，该语句等价于 $\forall x\,\atom{L}{x,x}$。现在，如果我们不够小心，忽略了 $a$ 确实出现在 $\forall x$ 的辖域内这一事实，那么 $\forall
x\,\atom{L}{x,x}$\emph{可以}通过将名称~$a$ 替换为变元~$x$ 从 $\forall
x\,\atom{L}{a,x}$ 得到。如果没有上述限制，那么下面的推导就会被允许：


\begin{fitchproof}
		\hypo{a}{\forall x\, \atom{L}{a,x}}\PR
		\ifHTMLtarget
		\have{e}{\exists x\, \forall x\, \atom{L}{x,x}}\by{不当的尝试：意图应用 $\exists$I}{a}
	\else
		\have{e}{\exists x\, \forall x\, \atom{L}{x,x}}\by{不当的尝试}{}
		\have[\ ]{x}{}\by{意图应用 $\exists$I}{a}
	\fi\end{fitchproof}


但是，由于 $\exists x\, \forall x\, \atom{L}{x,x}$ 等价于 $\forall x\, \atom{L}{x,x}$，这将是一个无效的推理。为了明白为什么，不妨把 $L$ 理解为“喜欢”：从“Alex 喜欢每个人”并不能推出“每个人都喜欢每个人”。

类似地，在\cref{s:forallIntro}中，我们说过，$\metav{A}(\ldots
\metav{x} \ldots \metav{x}\ldots)$ 是将 $\metav{A}(\ldots \metav{c}
\ldots \metav{c}\ldots)$ 中 $\metav{c}$ 的每一次出现替换为 $\metav{x}$ 的结果，但\emph{仅当 $\metav{c}$ 在 $\metav{A}$ 中不处于约束 $\metav{x}$ 的量词的辖域内时才允许这样做}。正如 $\exists$I 规则陈述中涉及的约束那样，在实践中你很少会面临违反它的风险。但这就说明了它为何重要。同样，由于我们允许空量化，$\forall
x\exists x\,\atom{L}{x,x}$ 是一阶逻辑的合法语句。然而，$\forall x\exists x\,\atom{L}{x,x}$ 可以通过将名称~$a$ 全部替换为变元~$x$ 从 $\forall
x\,\atom{L}{a,x}$ 得到。这违反了我们设定的限制：在这个例子中，$a$ 确实出现在 $\forall x$ 的辖域内。如果没有这个限制，下面的推导就会被允许：


\begin{fitchproof}
	\hypo{p}{\forall y\exists x\,\atom{L}{y,x}}\PR
	\have{a}{\exists x\, \atom{L}{a,x}}\Ae{p}
\ifHTMLtarget
	\have{e}{\forall x\, \exists x\, \atom{L}{x,x}}\by{不当的尝试：意图应用 $\forall$I}{a}
\else
\have{e}{\forall x\, \exists x\, \atom{L}{x,x}}\by{不当的尝试}{}
	\have[\ ]{x}{}\by{意图应用 $\forall$I}{a}
\fi
\end{fitchproof}


但是，由于 $\forall x\, \exists x\, \atom{L}{x,x}$ 等价于 $\exists x\, \atom{L}{x,x}$，这将是一个无效的推理：“某人喜欢自己”不能从“每个人都喜欢某个人”推出。

\practiceproblems
\problempart
解释为什么下面这两个“证明”是\emph{错误的}。同时，提供使得这些“证明”所体现的谬误论证形式无效的解释：


\begin{multicols}{2}
	\begin{fitchproof}
		\hypo{Rxx}{\forall x\, \atom{R}{x,x}}\PR
		\have{Raa}{\atom{R}{a,a}}\Ae{Rxx}
		\have{Ray}{\forall y\, \atom{R}{a,y}}\Ai{Raa}
		\have{Rxy}{\forall x\, \forall y\, \atom{R}{x,y}}\Ai{Ray}
	\end{fitchproof}
	\begin{fitchproof}
		\hypo{AE}{\forall x\, \exists y\, \atom{R}{x,y}}\PR
		\have{E}{\exists y\, \atom{R}{a,y}}\Ae{AE}
		\open
			\hypo{ass}{\atom{R}{a,a}}\AS
			\have{Ex}{\exists x\, \atom{R}{x,x}}\Ei{ass}
		\close
		\have{con}{\exists x\, \atom{R}{x,x}}\Ee{E, ass-Ex}
	\end{fitchproof}
\end{multicols}

\problempart
\label{pr.justifyFOLproof}
下面的三个证明缺少引用（规则和行号）。请补充它们，使其成为合规的证明。


\begin{compactlist}
\item \begin{fitchproof}
\hypo{p1}{\forall x\exists y(\atom{R}{x,y} \eor \atom{R}{y,x})}
\hypo{p2}{\forall x\,\enot \atom{R}{m,x}}
\have{3}{\exists y(\atom{R}{m,y} \eor \atom{R}{y,m})}{}
	\open
		\hypo{a1}{\atom{R}{m,a} \eor \atom{R}{a,m}}
		\have{a2}{\enot \atom{R}{m,a}}{}
		\have{a3}{\atom{R}{a,m}}{}
		\have{a4}{\exists x\, \atom{R}{x,m}}{}
	\close
\have{n}{\exists x\, \atom{R}{x,m}} {}
\end{fitchproof}

\item \begin{fitchproof}
\hypo{1}{\forall x(\exists y\,\atom{L}{x,y} \eif \forall z\,\atom{L}{z,x})}
\hypo{2}{\atom{L}{a,b}}
\have{3}{\exists y\,\atom{L}{a,y} \eif \forall z\atom{L}{z,a}}{}
\have{4}{\exists y\, \atom{L}{a,y}} {}
\have{5}{\forall z\, \atom{L}{z,a}} {}
\have{6}{\atom{L}{c,a}}{}
\have{7}{\exists y\,\atom{L}{c,y} \eif \forall z\,\atom{L}{z,c}}{}
\have{8}{\exists y\, \atom{L}{c,y}}{}
\have{9}{\forall z\, \atom{L}{z,c}}{}
\have{10}{\atom{L}{c,c}}{}
\have{11}{\forall x\, \atom{L}{x,x}}{}
\end{fitchproof}

\item \begin{fitchproof}
\hypo{a}{\forall x(\atom{J}{x} \eif \atom{K}{x})}
\hypo{b}{\exists x\,\forall y\, \atom{L}{x,y}}
\hypo{c}{\forall x\, \atom{J}{x}}
\open
	\hypo{2}{\forall y\, \atom{L}{a,y}}
	\have{3}{\atom{L}{a,a}}{}
	\have{d}{\atom{J}{a}}{}
	\have{e}{\atom{J}{a} \eif \atom{K}{a}}{}
	\have{f}{\atom{K}{a}}{}
	\have{4}{\atom{K}{a} \eand \atom{L}{a,a}}{}
	\have{5}{\exists x(\atom{K}{x} \eand \atom{L}{x,x})}{}
\close
\have{j}{\exists x(\atom{K}{x} \eand \atom{L}{x,x})}{}
\end{fitchproof}
\end{compactlist}

\problempart
\label{pr.BarbaraEtc.proof1}
在\hyperref[pr.BarbaraEtc]{练习 \ref*{s:MoreMonadic}A}中，我们考察了 Aristotle 逻辑的十五个三段论形式。请为每一种论证形式提供证明。注意：如果你将（例如）“没有F是G”符号化为$\forall x (\atom{F}{x} \eif \enot \atom{G}{x})$，你会发现证明要容易\emph{得多}。

\

\problempart
\label{pr.BarbaraEtc.proof2}
Aristotle 及其后继者们还识别出其他一些依赖于“存在涵义”的三段论形式。请用一阶逻辑符号化下述每一种论证形式，并给出证明。


\begin{compactlist}
	\item \textbf{Barbari.} 有H。所有G都是F。所有H都是G。所以：有H是F。
	\item \textbf{Celaront.} 有H。没有G是F。所有H都是G。所以：有H不是F。
	\item \textbf{Cesaro.} 有H。没有F是G。所有H都是G。所以：有H不是F。
	\item \textbf{Camestros.} 有H。所有F都是G。没有H是G。所以：有H不是F。
	\item \textbf{Felapton.} 有G。没有G是F。所有G都是H。所以：有H不是F。
	\item \textbf{Darapti.} 有G。所有G都是F。所有G都是H。所以：有H是F。
	\item \textbf{Calemos.} 有H。所有F都是G。没有G是H。所以：有H不是F。
	\item \textbf{Fesapo.} 有G。没有F是G。所有G都是H。所以：有H不是F。
	\item \textbf{Bamalip.} 有F。所有F都是G。所有G都是H。所以：有H是F。
\end{compactlist}

\problempart
\label{pr.someFOLproofs}
对于下列各断言，请提供一个一阶逻辑证明以表明其为真。


\begin{compactlist}
\item $\proves \forall x\,\atom{F}{x} \eif \forall y(\atom{F}{y} \eand \atom{F}{y})$
\item $\forall x(\atom{A}{x}\eif \atom{B}{x}), \exists x\,\atom{A}{x} \proves \exists x\,\atom{B}{x}$
\item $\forall x(\atom{M}{x} \eiff \atom{N}{x}), \atom{M}{a} \eand \exists x\,\atom{R}{x,a} \proves \exists x\,\atom{N}{x}$
\item $\forall x\, \forall y\,\atom{G}{x,y}\proves\exists x\,\atom{G}{x,x}$
\item $\proves\forall x\,\atom{R}{x,x} \eif \exists x\, \exists y\,\atom{R}{x,y}$
\item $\proves\forall y\, \exists x (\atom{Q}{y} \eif \atom{Q}{x})$
\item $\atom{N}{a} \eif \forall x(\atom{M}{x} \eiff \atom{M}{a}), \atom{M}{a}, \enot\atom{M}{b}\proves \enot \atom{N}{a}$
\item $\forall x\, \forall y (\atom{G}{x,y} \eif \atom{G}{y,x}) \proves \forall x\forall y (\atom{G}{x,y} \eiff \atom{G}{y,x})$
\item $\forall x(\enot\atom{M}{x} \eor \atom{L}{j,x}), \forall x(\atom{B}{x}\eif \atom{L}{j,x}), \forall x(\atom{M}{x}\eor \atom{B}{x})\proves \forall x\,\atom{L}{j,x}$
\end{compactlist}

\solutions
\problempart
\label{pr.likes}
为以下论证写一套符号化约定，将其符号化，并证明它：


\begin{earg}
\item 有这样一个人，她喜欢每一个喜欢所有她所喜欢的人的人。
\item[\texttherefore] 有这样一个人，她喜欢她自己。
\end{earg}

\problempart
证明下列各对语句是可互推的。


\begin{compactlist}
\item $\forall x (\atom{A}{x}\eif \enot \atom{B}{x})$, $\enot\exists x(\atom{A}{x} \eand \atom{B}{x})$
\item $\forall x (\enot\atom{A}{x}\eif \atom{B}{d})$, $\forall x\,\atom{A}{x} \eor \atom{B}{d}$
\item $\exists x\,\atom{P}{x} \eif \atom{Q}{c}$, $\forall x (\atom{P}{x} \eif \atom{Q}{c})$
\end{compactlist}

\solutions
\problempart
\label{pr.FOLequivornot}
对于下列各对语句：如果它们可互推，请给出证明以表明这一点。如果它们不可互推，请构造一个解释以表明它们在逻辑上不等价。


\begin{compactlist}
\item $\forall x\,\atom{P}{x} \eif \atom{Q}{c}, \forall x (\atom{P}{x} \eif \atom{Q}{c})$
\item $\forall x\,\forall y\, \forall z\,\atom{B}{x,y,z}, \forall x\,\atom{B}{x,x,x}$
\item $\forall x\,\forall y\,\atom{D}{x,y}, \forall y\,\forall x\,\atom{D}{x,y}$
\item $\exists x\,\forall y\,\atom{D}{x,y}, \forall y\,\exists x\,\atom{D}{x,y}$
\item $\forall x (\atom{R}{c,a} \eiff \atom{R}{x,a}), \atom{R}{c,a} \eiff \forall x\,\atom{R}{x,a}$
\end{compactlist}

\solutions
\problempart
\label{pr.FOLvalidornot}
对于下列各论证：如果它在一阶逻辑中有效，请给出证明。如果它无效，请构造一个解释以表明其无效。


\begin{compactlist}
\item $\exists y\,\forall x\,\atom{R}{x,y} \therefore \forall x\,\exists y\,\atom{R}{x,y}$
\item $\forall x\,\exists y\,\atom{R}{x,y} \therefore  \exists y\,\forall x\,\atom{R}{x,y}$
\item $\exists x(\atom{P}{x} \eand \enot \atom{Q}{x}) \therefore \forall x(\atom{P}{x} \eif \enot \atom{Q}{x})$
\item $\forall x(\atom{S}{x} \eif \atom{T}{a}), \atom{S}{d} \therefore \atom{T}{a}$
\item $\forall x(\atom{A}{x}\eif \atom{B}{x}), \forall x(\atom{B}{x} \eif \atom{C}{x}) \therefore \forall x(\atom{A}{x} \eif \atom{C}{x})$
\item $\exists x(\atom{D}{x} \eor \atom{E}{x}), \forall x(\atom{D}{x} \eif \atom{F}{x}) \therefore \exists x(\atom{D}{x} \eand \atom{F}{x})$
\item $\forall x\,\forall y(\atom{R}{x,y} \eor \atom{R}{y,x}) \therefore \atom{R}{j,j}$
\item $\exists x\,\exists y(\atom{R}{x,y} \eor \atom{R}{y,x}) \therefore \atom{R}{j,j}$
\item $\forall x\,\atom{P}{x} \eif \forall x\,\atom{Q}{x}, \exists x\, \enot\atom{P}{x} \therefore \exists x\, \enot \atom{Q}{x}$
\item $\exists x\,\atom{M}{x} \eif \exists x\,\atom{N}{x}$, $\enot \exists x\,\atom{N}{x}\therefore  \forall x\, \enot \atom{M}{x}$
\end{compactlist}

\chapter{带量词的证明}

在\cref{s:stratTFL}中，我们讨论了使用真值函项逻辑的自然演绎基本规则构造证明的策略。同样的原则也适用于量词的规则。如果我们想证明一个量化句子$\forall \metav{x}\,
\atom{\metav{A}}{\metav{x}}$或$\exists \metav{x}\,
\atom{\metav{A}}{\metav{x}}$，我们可以通过$\forall$I或$\exists$I来证成我们想要的句子，并尝试找到该规则对应前提的证明，以此进行逆向推理。而从量化句子进行正向推理时，我们则根据情况应用$\forall$E或$\exists$E。

具体来说，假设你想证明$\forall \metav{x}\, \atom{\metav{A}}{\metav{x}}$。要使用$\forall$I来证明，我们需要对某个名称~$\metav{c}$证明$\atom{\metav{A}}{\metav{c}}$，该名称不能出现在任何未解除的假设中。应用相应的策略，即通过逆向推理构造$\forall \metav{x}\, \atom{\metav{A}}{\metav{x}}$的证明，就是在它上面写下$\atom{\metav{A}}{\metav{c}}$，然后继续尝试找到该句子的证明。


\begin{fitchproof}
	\ellipsesline
	\have[n]{n}{\atom{\metav{A}}{\metav{c}}}
	\have{m}{\forall \metav{x}\, \atom{\metav{A}}{\metav{x}}}\Ai{n}
\end{fitchproof}


$\atom{\metav{A}}{\metav{c}}$是通过将$\atom{\metav{A}}{\metav{x}}$中$\metav{x}$的每个自由出现替换为~$\metav{c}$而从$\atom{\metav{A}}{\metav{x}}$得到的。要使这可行，$\metav{c}$必须满足特殊条件。我们可以通过总是选择一个尚未在目前已构造的证明中出现的名称来确保这一点。（当然，它最终会出现在我们构造的证明中——只是不出现在第~$n+1$行未解除的假设中。）

要从句子$\exists \metav{x}\, \atom{\metav{A}}{\metav{x}}$逆向推理，我们类似地在它上面写下一个句子，该句子可以作为应用$\exists$I规则的证成依据，即形式为$\atom{\metav{A}}{\metav{c}}$的句子。


\begin{fitchproof}
	\ellipsesline
	\have[n]{n}{\atom{\metav{A}}{\metav{c}}}
	\have{m}{\exists \metav{x}\, \atom{\metav{A}}{\metav{x}}}\Ei{n}
\end{fitchproof}


这看起来就像我们从全称量化句子逆向推理时所作的一样。区别在于，对于$\forall$I，我们必须选择一个尚未在证明中（到目前为止）出现的名称~$\metav{c}$；而对于$\exists$I，我们可以且通常必须选择一个已经在证明中出现的名称~$\metav{c}$。就像在$\eor$I的情况下一样，通常不清楚哪个$\metav{c}$会奏效，因此为了避免回溯，你应该只在所有其他策略都已应用后才从存在量化句子逆向推理。

相比之下，从句子$\exists \metav{x}\, \atom{\metav{A}}{\metav{x}}$\emph{正向推理}通常总是有效的，你不需要回溯。从存在量化句子正向推理不仅考虑$\exists \metav{x}\, \atom{\metav{A}}{\metav{x}}$，还考虑你希望证明的任何句子$\metav{B}$。它要求你在$\metav{B}$之上建立一个子证明，其中$\metav{B}$是最后一行，并且以$\exists \metav{x}\, \atom{\metav{A}}{\metav{x}}$的一个代入实例$\atom{\metav{A}}{\metav{c}}$作为假设。为了确保$\exists$E规则对$\metav{c}$所施加的条件得到满足，请选择一个尚未在证明中出现的名称$\metav{c}$。


\begin{fitchproof}
	\ellipsesline
	\have[m]{m}{\exists \metav{x}\, \atom{\metav{A}}{\metav{x}}}
	\ellipsesline
	\open
	\hypo[n]{n}{\atom{\metav{A}}{\metav{c}}}\AS
	\ellipsesline
	\have[k]{k}{\metav{B}}
	\close
	\have{e}{\metav{B}}\Ee{m,n-k}
\end{fitchproof}


然后你将继续以证明$\metav{B}$为目标，但现在是在一个子证明中，你有一个额外的句子可以使用，即~$\atom{\metav{A}}{\metav{c}}$。

最后，从$\forall \metav{x}\, \atom{\metav{A}}{\metav{x}}$正向推理意味着对于任何名称~$\metav{c}$，你总是可以写下$\atom{\metav{A}}{\metav{c}}$并用$\forall$E证成它。当然，你不会想不加选择地这样做。只有某些名称$\metav{c}$会帮助你证明你正在处理的目标句子。因此，就像从$\exists \metav{x}\, \atom{\metav{A}}{\metav{x}}$逆向推理一样，你应该只在所有其他策略都已应用后才从$\forall \metav{x}\, \atom{\metav{A}}{\metav{x}}$正向推理。

让我们以论证$\forall x(\atom{A}{x} \eif B) \therefore \exists x\,\atom{A}{x} \eif B$为例。开始构造证明时，我们将前提写在顶部，结论写在底部。


\begin{fitchproof}
\hypo{1}{\forall x(\atom{A}{x} \eif B)}\PR
\ellipsesline
\have[n]{7}{\exists x\,\atom{A}{x} \eif B}
\end{fitchproof}


真值函项逻辑联结词的策略仍然适用，你应该以同样的顺序应用它们：首先从条件句、否定句、合取句以及现在还有全称量词逆向推理，然后从析取句以及现在存在量词正向推理，只有在这之后才尝试应用$\eif$E、$\enot$E、$\lor$I、$\forall$E或$\exists$I。在我们的例子中，这意味着从结论逆向推理：


\begin{fitchproof}
	\hypo{1}{\forall x(\atom{A}{x} \eif B)}\PR
	\open
	\hypo{2}{\exists x\,\atom{A}{x}}\AS
	\ellipsesline
	\have[n][-1]{6}{B}
	\close
	\have[n]{7}{\exists x\,\atom{A}{x} \eif B}\ci{2-(6)}
\end{fitchproof}


我们接下来的步骤应该从第$2$行的$\exists x\,\atom{A}{x}$正向推理。为此，我们必须选择一个尚未出现在证明中的名称。既然还没有任何名称出现，我们可以选任意一个名称，比方说$d$：


\begin{fitchproof}
	\hypo{1}{\forall x(\atom{A}{x} \eif B)}\PR
	\open
	\hypo{2}{\exists x\,\atom{A}{x}}\AS
	\open
	\hypo{3}{\atom{A}{d}}\AS
	\ellipsesline
	\have[n][-2]{5}{B}
	\close
	\have[n][-1]{6}{B}\Ee{2,3-(5)}
	\close
	\have[n]{7}{\exists x\,\atom{A}{x} \eif B}\ci{2-(6)}
\end{fitchproof}


现在我们已经用完了主要策略，是时候从前提$\forall x(\atom{A}{x} \eif B)$正向推理了。应用$\forall$E意味着我们可以为$A(\metav{c}) \eif B$的任何实例提供证成，无论我们选择哪个$\metav{c}$。当然，明智的做法是选择$d$，因为这将给我们$\atom{A}{d} \eif B$。接着我们就可以应用$\eif$E来证成$B$，从而完成证明。


\begin{fitchproof}
	\hypo{1}{\forall x(\atom{A}{x} \eif B)}\PR
	\open
	\hypo{2}{\exists x\,\atom{A}{x}}\AS
	\open
	\hypo{3}{\atom{A}{d}}\AS
\have{4}{\atom{A}{d} \eif B}\Ae{1}
	\have{5}{B}\ce{4,3}
	\close
	\have{6}{B}\Ee{2,3-5}
	\close
	\have{7}{\exists x\,\atom{A}{x} \eif B}\ci{2-6}
\end{fitchproof}

现在让我们来构造其逆命题的证明。我们从以下形式开始：


\begin{fitchproof}
	\hypo{1}{\exists x\,\atom{A}{x} \eif B}\PR
	\ellipsesline
	\have[n]{7}{\forall x(\atom{A}{x} \eif B)}
\end{fitchproof}


注意，前提是一个条件句，而不是一个存在量化句子，因此我们（暂时）还不应该从它正向推理。从结论$\forall x(\atom{A}{x} \eif B)$逆向推理，引导我们去寻找一个对$\atom{A}{d} \eif B$的证明：


\begin{fitchproof}
	\hypo{1}{\exists x\,\atom{A}{x} \eif B}\PR
	\ellipsesline
	\have[n][-1]{6}{\atom{A}{d} \eif B}
	\have[n]{7}{\forall x(\atom{A}{x} \eif B)}\Ai{6}
\end{fitchproof}


而从$\atom{A}{d} \eif B$逆向推理，意味着我们应该建立一个以$\atom{A}{d}$为假设、以$B$为最后一行的子证明：


\begin{fitchproof}
	\hypo{1}{\exists x\,\atom{A}{x} \eif B}\PR
	\open
	\hypo{2}{\atom{A}{d}}\AS
	\ellipsesline
	\have[n][-2]{5}{B}
	\close
	\have[n][-1]{6}{\atom{A}{d} \eif B}\ci{2-(5)}
	\have[n]{7}{\forall x(\atom{A}{x} \eif B)}\Ai{6}
\end{fitchproof}


现在我们可以从第$1$行的前提正向推理了。那是一个条件句，而它的后件恰好就是我们想要证成的句子$B$。所以我们应该寻找一个对其前件$\exists x\,\atom{A}{x}$的证明。当然，这现在已是现成可用的，对第$2$行应用$\exists$I即可得到。这样我们就完成了：


\begin{fitchproof}
	\hypo{1}{\exists x\,\atom{A}{x} \eif B}\PR
	\open
	\hypo{2}{\atom{A}{d}}\AS
	\have{3}{\exists x\,\atom{A}{x}}\Ei{2}
	\have{5}{B}\ce{1,3}
	\close
	\have{6}{\atom{A}{d} \eif B}\ci{2-5}
	\have{7}{\forall x(\atom{A}{x} \eif B)}\Ai{6}
\end{fitchproof}

\practiceproblems

\problempart
运用策略为以下每一个论证和定理寻找证明：


\begin{compactlist}
\item $A \eif \forall x\,\atom{B}{x} \therefore \forall x(A \eif \atom{B}{x})$
\item $\exists x(A \eif \atom{B}{x}) \therefore A \eif \exists x\, \atom{B}{x}$
\item $\forall x(\atom{A}{x} \eand \atom{B}{x}) \eiff (\forall x\,\atom{A}{x} \eand \forall x\,\atom{B}{x})$
\item $\exists x(\atom{A}{x} \eor \atom{B}{x}) \eiff (\exists x\,\atom{A}{x} \eor \exists x\,\atom{B}{x})$
\item $A \eor \forall x\,\atom{B}{x}) \therefore \forall x(A \eor \atom{B}{x})$
\item $\forall x(\atom{A}{x} \eif B) \therefore \exists x\,\atom{A}{x} \eif B$
\item $\exists x(\atom{A}{x} \eif B) \therefore \forall x\,\atom{A}{x} \eif B$
\item $\forall x(\atom{A}{x} \eif \exists y\,\atom{A}{y})$
\end{compactlist}


除基本量词规则外，仅使用真值函项逻辑的基本规则。

\problempart
运用策略为以下每一个论证和定理寻找证明：


\begin{compactlist}
\item $\forall x\,\atom{R}{x,x} \therefore \forall x\,\exists y\,\atom{R}{x,y}$
\item $\forall x\,\forall y\,\forall z[(\atom{R}{x,y} \eand \atom{R}{y,z}) \eif \atom{R}{x,z}]$ \\
$\therefore \forall x\,\forall y[\atom{R}{x,y} \eif \forall z(\atom{R}{y,z} \eif \atom{R}{x,z})]$
\item $\forall x\,\forall y\,\forall z[(\atom{R}{x,y} \eand \atom{R}{y,z}) \eif \atom{R}{x,z}],$\\ $\forall x\,\forall y(\atom{R}{x,y} \eif \atom{R}{y, x})$ \\ $\therefore \forall x\,\forall y\,\forall z[(\atom{R}{x,y} \eand \atom{R}{x,z}) \eif \atom{R}{y,z}]$
\item $\forall x\,\forall y(\atom{R}{x,y} \eif \atom{R}{y, x})$ \\$\therefore \forall x\,\forall y\,\forall z[(\atom{R}{x,y} \eand \atom{R}{x,z}) \eif \exists u(\atom{R}{y,u} \eand \atom{R}{z,u})]$
\item $\enot \exists x\,\forall y (\atom{A}{x, y} \eiff \lnot\atom{A}{y, y})$
\end{compactlist}

\problempart
运用策略为以下每一个论证和定理寻找证明：


\begin{compactlist}
\item $\forall x\,\atom{A}{x} \eif B \therefore \exists x(\atom{A}{x} \eif B)$
\item $A \eif \exists x\, \atom{B}{x} \therefore \exists x(A \eif \atom{B}{x})$
\item $\forall x(A \eor \atom{B}{x}) \therefore A \eor \forall x\,\atom{B}{x})$
\item $\exists x(\atom{A}{x} \eif \forall y\,\atom{A}{y})$
\item $\exists x(\exists y\,\atom{A}{y} \eif \atom{A}{x})$
\end{compactlist}


这些证明需要使用间接证明法。除基本量词规则外，仅使用真值函项逻辑的基本规则。

\chapter{量词转换}\label{s:CQ}

在这一节，我们将在前一节基本规则的基础上增加一些额外规则。这些规则规范量词和否定之间的相互作用。

在\cref{s:FOLBuildingBlocks}中，我们注意到$\enot\exists x\metav{A}$逻辑等价于$\forall x\, \enot\metav{A}$。我们将在证明系统中加入一些规则来规范这一点。具体来说，我们增加：
	\factoidbox{


\begin{fitchproof}
		\have[m]{a}{\forall \metav{x}\, \enot\metav{A}}
		\have[\ ]{con}{\enot \exists \metav{x}\, \metav{A}}\cq{a}
	\end{fitchproof}

}
和
\factoidbox{


\begin{fitchproof}
		\have[m]{a}{ \enot \exists \metav{x}\, \metav{A}}
		\have[\ ]{con}{\forall \metav{x}\, \enot \metav{A}}\cq{a}
	\end{fitchproof}

}
同样，我们增加：
\factoidbox{


\begin{fitchproof}
		\have[m]{a}{\exists \metav{x}\, \enot \metav{A}}
		\have[\ ]{con}{\enot \forall \metav{x}\, \metav{A}}\cq{a}
	\end{fitchproof}

}
和
\factoidbox{


\begin{fitchproof}
		\have[m]{a}{\enot \forall \metav{x}\, \metav{A}}
		\have[\ ]{con}{\exists \metav{x}\, \enot \metav{A}}\cq{a}
	\end{fitchproof}

}
\practiceproblems
\problempart
证明下列各组句子是不一致的：

\begin{compactlist}
\item $\atom{S}{a}\eif \atom{T}{m}, \atom{T}{m} \eif \atom{S}{a}, \atom{T}{m} \eand \enot \atom{S}{a}$
\item $\enot\exists x\,\atom{R}{x,a}, \forall x\, \forall y\,\atom{R}{y,x}$
\item $\enot\exists x\, \exists y\,\atom{L}{x,y}, \atom{L}{a,a}$
\item $\forall x(\atom{P}{x} \eif \atom{Q}{x}), \forall z(\atom{P}{z} \eif \atom{R}{z}), \forall y\,\atom{P}{y}, \enot \atom{Q}{a} \eand \enot \atom{R}{b}$
\end{compactlist}

\problempart
证明下列每对句子是可互推的：

\begin{compactlist}
\item $\forall x (\atom{A}{x}\eif \enot \atom{B}{x}), \enot\exists x(\atom{A}{x} \eand \atom{B}{x})$
\item $\forall x (\enot\atom{A}{x}\eif \atom{B}{d}), \forall x\,\atom{A}{x} \eor \atom{B}{d}$
\end{compactlist}

\problempart
在\cref{s:MoreMonadic}中，我们考虑了将量词“移过”各种逻辑算子会发生什么。证明下列每对句子是可互推的：

\begin{compactlist}
\item $\forall x(\atom{F}{x} \eand \atom{G}{a}), \forall x\,\atom{F}{x} \eand \atom{G}{a}$
\item $\exists x(\atom{F}{x} \eor \atom{G}{a}), \exists x\,\atom{F}{x} \eor \atom{G}{a}$
\item $\forall x(\atom{G}{a} \eif \atom{F}{x}), \atom{G}{a} \eif \forall x\,\atom{F}{x}$
\item $\forall x(\atom{F}{x} \eif \atom{G}{a}), \exists x\,\atom{F}{x} \eif \atom{G}{a}$
\item $\exists x(\atom{G}{a} \eif \atom{F}{x}), \atom{G}{a} \eif \exists x\,\atom{F}{x}$
\item $\exists x(\atom{F}{x} \eif \atom{G}{a}), \forall x\,\atom{F}{x} \eif \atom{G}{a}$
\end{compactlist}

注意：变元 $x$ 在 $\atom{G}{a}$ 中不出现。当一个句子的所有量词都出现在句首时，这个句子被称为处于\emph{前束范式}。这些等价式有时也被称为\emph{前束化规则}，因为它们为我们提供了将任何句子化为前束范式的方法。

\chapter{同一性规则}
在\cref{s:Interpretations}中，我们提到了一个在哲学上有争议的论点，即\emph{不可分辨者的同一性}。这个论点声称，在所有方面都不可分辨的对象实际上是彼此同一的。我们也提到，我们不会采纳这个论点。因此，无论你了解到关于两个对象的多少信息，我们都不能证明它们是同一的。当然，除非你本来就已知道这两个对象实际上是同一的，但那样的话，证明也就没什么启发性了。

不过，总的要点是，\emph{没有任何}不包含同一性谓词的句子能够为推断出 $a=b$ 提供理据。所以，我们的同一性引入规则不能允许我们推断出一个包含两个\emph{不同}名称的同一性陈述。

然而，每个对象都与它自身同一。要得出某物与自身同一的结论，并不需要任何前提。因此，这将是我们的同一性引入规则：
\factoidbox{


\begin{fitchproof}
	\have[\ \,\,\,]{x}{\metav{c}=\metav{c}} \by{=I}{}
\end{fitchproof}

}
注意，这条规则不要求引用证明中任何先前的行。对于任何名称 \metav{c}，你都可以在任何一行上写下 $\metav{c}=\metav{c}$，并以 {=}I 规则作为唯一的理据。

我们的消去规则则更有意思。如果你已确立了 $a=b$，那么任何对于以 $a$ 命名的对象为真的东西，也必定对以 $b$ 命名的对象为真。对于任何一个包含 $a$ 的句子，你可以把其中一些或所有出现的 $a$ 替换为 $b$，从而得到一个等价的句子。例如，从 $\atom{R}{a,a}$ 和 $a = b$ 出发，你有理据推断出 $\atom{R}{a,b}$、$\atom{R}{b,a}$ 或 $\atom{R}{b,b}$。更一般地：
\factoidbox{

\begin{fitchproof}
	\have[m]{e}{\metav{a}=\metav{b}}
	\have[n]{a}{\metav{A}(\ldots \metav{a} \ldots \metav{a}\ldots)}
	\have[\ ]{ea1}{\metav{A}(\ldots \metav{b} \ldots \metav{a}\ldots)} \by{=E}{e,a}
\end{fitchproof}

}
这里的记法与 $\exists$I 规则中的类似。因此，$\metav{A}(\ldots \metav{a} \ldots \metav{a}\ldots)$ 是包含名称 $\metav{a}$ 的公式，而 $\metav{A}(\ldots \metav{b} \ldots \metav{a}\ldots)$ 则是通过将该公式中的一个或多个 $\metav{a}$ 实例替换为名称 $\metav{b}$ 而得到的公式。第 $m$ 行和第 $n$ 行出现的顺序不限，也不需相邻，但我们引用时总是先引用同一性陈述。对称地，我们允许：
\factoidbox{

\begin{fitchproof}
	\have[m]{e}{\metav{a}=\metav{b}}
	\have[n]{a}{\metav{A}(\ldots \metav{b} \ldots \metav{b}\ldots)}
	\have[\ ]{ea2}{\metav{A}(\ldots \metav{a} \ldots \metav{b}\ldots)} \by{=E}{e,a}
\end{fitchproof}

}
这条规则有时被称为\emph{Leibniz 律}，以 Gottfried Leibniz 命名。

为了展示这些规则的作用，我们将证明一些简单的结果。首先，我们证明同一性是\emph{对称的}：


\begin{fitchproof}
	\open
		\hypo{ab}{a = b}\AS
		\have{aa}{a = a}\by{=I}{}
		\have{ba}{b = a}\by{=E}{ab, aa}
	\close
	\have{abba}{a = b \eif b =a}\ci{ab-ba}
	\have{ayya}{\forall y (a = y \eif y = a)}\Ai{abba}
	\have{xyyx}{\forall x\, \forall y (x = y \eif y = x)}\Ai{ayya}
\end{fitchproof}


我们通过将第 $2$ 行中的一个 $a$ 实例替换为一个 $b$ 实例而得到第 $3$ 行；这在给定 $a= b$ 的情况下是合理的。

其次，我们证明同一性是\emph{传递的}：


\begin{fitchproof}
	\open
		\hypo{abc}{a = b \eand b = c}\AS
		\have{ab}{a = b}\ae{abc}
		\have{bc}{b = c}\ae{abc}
		\have{ac}{a = c}\by{=E}{ab, bc}
	\close
	\have{con}{(a = b \eand b =c) \eif a = c}\ci{abc-ac}
	\have{conz}{\forall z((a = b \eand b = z) \eif a = z)}\Ai{con}
	\have{cony}{\forall y\,\forall z((a = y \eand y = z) \eif a = z)}\Ai{conz}
	\have{conx}{\forall x\,\forall y \forall z((x = y \eand y = z) \eif x = z)}\Ai{cony}
\end{fitchproof}


我们通过将第 $3$ 行中的 $b$ 替换为 $a$ 而得到第 $4$ 行；这在给定 $a= b$ 的情况下是合理的。

\practiceproblems
\problempart
\label{pr.identity}
针对下列每个断言，提供一个显示其为真的一阶逻辑证明。
\begin{compactlist}
\item $\atom{P}{a} \eor \atom{Q}{b}, \atom{Q}{b} \eif b=c, \enot\atom{P}{a} \proves \atom{Q}{c}$
\item $m=n \eor n=o, \atom{A}{n} \proves \atom{A}{m} \eor \atom{A}{o}$
\item $\forall x\ x=m, \atom{R}{m,a} \proves \exists x\,\atom{R}{x,x}$
\item $\forall x\,\forall y(\atom{R}{x,y} \eif x=y)\proves \atom{R}{a,b} \eif \atom{R}{b,a}$
\item $\enot \exists x\enot\, {x = m} \proves \forall x\,\forall y (\atom{P}{x} \eif \atom{P}{y})$
\item $\exists x\,\atom{J}{x}, \exists x\, \enot\atom{J}{x}\proves \exists x\, \exists y\, \enot\, {x = y}$
\item $\forall x(x=n \eiff \atom{M}{x}), \forall x(\atom{O}{x} \eor \enot \atom{M}{x})\proves \atom{O}{n}$
\item $\exists x\,\atom{D}{x}, \forall x(x=p \eiff \atom{D}{x})\proves \atom{D}{p}$
\item $\exists x\bigl[(\atom{K}{x} \eand \forall y(\atom{K}{y} \eif x=y)) \eand \atom{B}{x}\bigr], \atom{K}{d} \proves \atom{B}{d}$
\item $\proves \atom{P}{a} \eif \forall x(\atom{P}{x} \eor \enot\, {x = a})$
\end{compactlist}

\problempart
证明下列句子是可互推的：
\begin{itemize}
\item $\exists x \bigl([\atom{F}{x} \eand \forall y (\atom{F}{y} \eif x = y)] \eand x = n\bigr)$
\item $\atom{F}{n} \eand \forall y (\atom{F}{y} \eif n= y)$
\end{itemize}
从而说明两者都有充分的理由可以符号化自然语言句子“Nick 是那个～$F$”。\\

\problempart
在\cref{sec.identity}中，我们声称下列都是自然语言句子“恰好存在一个 $F$”的逻辑等值符号化：
\begin{itemize}
\item $\exists x\,\atom{F}{x} \eand \forall x\, \forall y \bigl[(\atom{F}{x} \eand \atom{F}{y}) \eif x = y\bigr]$
\item $\exists x \bigl[\atom{F}{x} \eand \forall y (\atom{F}{y} \eif x = y)\bigr]$
\item $\exists x\, \forall y (\atom{F}{y} \eiff x = y)$
\end{itemize}
证明它们都是可互推的。（\emph{提示}：要证明三个断言可互推，只需证明第一个能推出第二个、第二个能推出第三个、第三个能推出第一个即可；想想这是为什么。）

\
\problempart
将以下论证符号化：
\begin{earg}
		\item 恰好存在一个 $F$。
		\item 恰好存在一个 $G$。
		\item 没有东西既是 $F$ 又是 $G$。
		\item[\texttherefore] 恰好存在两个东西是 $F$ 或 $G$。
	\end{earg}
并为其提供一个证明。
%\begin{itemize}
%\item  $\exists x \bigl[\atom{F}{x} \eand \forall y (\atom{F}{y} \eif x = y)\bigr], \exists x \bigl[\atom{G}{x} \eand \forall y (\atom{G}{y} \eif x = y)\bigr], \forall x (\enot\atom{F}{x} \eor \enot \atom{G}{x}) \proves \exists x\, \exists y \bigl[\enot\, {x = y} \eand \forall z ((\atom{F}{z} \eor \atom{G}{z}) \eif (x = y \eor x = z))\bigr]$
%\end{itemize}

\chapter{派生规则}\label{s:DerivedFOL}
与真值函项逻辑的情况一样，我们首先将一阶逻辑的某些规则作为基本规则引入（在\cref{s:BasicFOL}中），然后又为量词转换添加了一些进一步的规则（在\cref{s:CQ}中）。事实上，量词转换规则应被视为\emph{派生}规则，因为它们可以从\cref{s:BasicFOL}的\emph{基本}规则中推导出来。（这里的要点与\cref{s:Derived}相同。）以下是第一条量词转换规则的论证：
\begin{fitchproof}
	\have[m]{An}{\forall x\, \enot \atom{A}{x}}
	\open
		\hypo[m+1]{E}{\exists x\, \atom{A}{x}}\AS
		\open
			\hypo[m+2]{c}{\atom{A}{c}}\AS
			\have[m+3]{nc}{\enot \atom{A}{c}}\Ae{An}
			\have[m+4]{red}{\ered}\ne{nc,c}
		\close
		\have[m+5]{red2}{\ered}\Ee{E,(c)-(red)}
	\close
	\have[m+6]{dada}{\enot \exists x\, \atom{A}{x}}\ni{(E)-(red2)}
\end{fitchproof}
%You will note that on line 3 I have written `for $\exists$E'. This is not technically a part of the proof. It is just a reminder---to me and to you---of why I have bothered to introduce `$\enot \atom{A}{c}$' out of the blue. You might find it helpful to add similar annotations to assumptions when performing proofs. But do not add annotations on lines other than assumptions: the proof requires its own citation, and your annotations will clutter it.
以下是第三条量词转换规则的论证：
\begin{fitchproof}
	\have[m]{nEna}{\exists x\, \enot \atom{A}{x}}
	\open
		\hypo[m+1]{Aa}{\forall x\, \atom{A}{x}}\AS
		\open
			\hypo[m+2]{nac}{\enot \atom{A}{c}}\AS
			\have[m+3]{a}{\atom{A}{c}}\Ae{Aa}
			\have[m+4]{con}{\ered}\ne{nac,a}
		\close
		\have[m+5]{con1}{\ered}\Ee{nEna, (nac)-(con)}
	\close
	\have[m+6]{dada}{\enot \forall x\, \atom{A}{x}}\ni{(Aa)-(con1)}
\end{fitchproof}
这就解释了为什么量词转换规则可以被当作派生规则来对待。对于另外两条量词转换规则，也可以给出类似的论证。

\practiceproblems
\problempart
提供证明，以论证将第二条和第四条量词转换规则作为派生规则加入的合理性。

\chapter{证明与语义}
我们在本书中使用了两种不同的断定符。这个：
$$\metav{A}_1, \metav{A}_2, \ldots, \metav{A}_n \proves \metav{C}$$
意味着存在一个以 $\metav{C}$ 结尾的证明，且其所有未被解除的假设都在$\metav{A}_1, \metav{A}_2, \ldots, \metav{A}_n$之中。这是一个\emph{证明论概念}。与此相对，这个：
$$\metav{A}_1, \metav{A}_2, \ldots, \metav{A}_n \entails \metav{C}$$
意味着没有任何一个赋值（或解释）能使$\metav{A}_1, \metav{A}_2, \ldots, \metav{A}_n$全部为真而 $\metav{C}$ 为假。这关乎对句子的真值指派。它是一个\emph{语义概念}。

无论怎样强调这两个概念是不同的，都不过分。但我们还可以再强调一下：\emph{它们是不同的概念。}

一旦你内化了这一点，请继续往下读。

尽管我们的语义概念和证明论概念是不同的，但两者之间存在着深刻的联系。为了解释这种联系，我们将从思考有效式与定理之间的关系开始。

要证明一个句子是定理，你只需要构造一个证明。诚然，构造一个二十行的证明可能很难，但检查证明的每一行并确认其合法性却没那么难；如果证明的每一行本身都是合法的，那么整个证明就是合法的。然而，要证明一个句子是有效式，则需要考虑所有可能的解释。在证明一个句子是定理和证明它是有效式之间做选择，证明它是定理会更容易一些。

反过来，要证明一个句子\emph{不是}定理是很难的。我们需要考虑所有（可能的）证明。这是非常困难的。然而，要证明一个句子不是有效式，你只需要构造一个使该句子为假的解释。诚然，想出这个解释可能很难；但一旦你做到了，要检查它给句子指派了什么真值就相对直接了。在证明一个句子不是定理和证明它不是有效式之间做选择，证明它不是有效式会更容易一些。

幸运的是，\emph{一个语句是定理当且仅当它是有效式}。因此，如果我们不依赖任何假设给出了 $\metav{A}$ 的证明，从而表明 $\metav{A}$ 是定理，即 ${}\proves \metav{A}$，我们就可以合理地推断 $\metav{A}$ 是有效式，即 $\entails\metav{A}$。类似地，如果我们构造了一个使 \metav{A} 为假的解释，从而表明它不是有效式，即 $\nentails \metav{A}$，即可得出 \metav{A} 不是定理，即 $\nproves \metav{A}$。

更一般地说，我们有以下强有力的结果：
\factoidbox{$\metav{A}_1, \metav{A}_2, \ldots, \metav{A}_n
\proves\metav{B}$ \textbf{当且仅当} $\metav{A}_1, \metav{A}_2, \ldots,
\metav{A}_n \entails\metav{B}$}
这表明，尽管可证性和语义后承是\emph{不同的}概念，但它们在外延上是等价的。因此：
	\factoidbox{

\begin{factoiditems}
		\item 一个论证是\emph{有效的}当且仅当\emph{结论可以从前提证明出来}。
		\item 一个语句是\emph{有效式}当且仅当它是一个\emph{定理}。
		\item 两个语句是\emph{等价的}当且仅当它们是\emph{可互推的}。
		\item 一组语句是\emph{共同可满足的}当且仅当它们是\emph{共同一致的}。
	\end{factoiditems}

}
基于这个原因，你可以根据需要选择何时从证明的角度思考，何时从赋值/解释的角度思考，选择对给定任务更容易的角度。下一页的表格总结了哪种方法（通常）更容易。

可证性和语义后承应该一致，这在直觉上是合理的。但是——让我们再强调一遍——不要被符号$\entails$和$\proves$的相似性所蒙蔽。这两个符号的含义非常不同。可证性和语义后承一致这一事实并非唾手可得的结果。

事实上，证明可证性与语义后承一致，恰恰是入门逻辑转变为中级逻辑的决定性标志。

\ifHTMLtarget


\begin{table}
\else
\begin{sidewaystable}\small
\fi
	\begin{center}
\begin{tabular*}{\textwidth}{p{.25\textheight}p{.325\textheight}p{.325\textheight}}
 & \textbf{是}  & \textbf{否}\\
\\
\metav{A} 是\textbf{有效式}吗？
& 给出一个表明 $\proves\metav{A}$ 的证明
& 给出一个使得 \metav{A} 为假的解释\\
\\
\metav{A} 是\textbf{矛盾式}吗？ &
给出一个表明 $\metav{A} \proves \ered$ 的证明 &
给出一个使得 \metav{A} 为真的解释\\
\\
%Is \metav{A} contingent? &
%give two interpretations, one in which \metav{A} is true and another in which \metav{A} is false & give a proof which either shows $\proves\metav{A}$ or $\proves\enot\metav{A}$\\
%\\
\metav{A} 和 \metav{B} 是\textbf{等价的}吗？ &
给出两个证明，一个证明 $\metav{A}\proves\metav{B}$，另一个证明 $\metav{B}\proves\metav{A}$
& 给出一个使得 \metav{A} 和 \metav{B} 真值不同的解释\\
\\
$\metav{A}_1, \metav{A}_2, \ldots, \metav{A}_n$ 是\textbf{共同可满足的}吗？
& 给出一个使得所有 $\metav{A}_1, \metav{A}_2, \ldots, \metav{A}_n$ 都为真的解释
& 从假设 $\metav{A}_1, \metav{A}_2, \ldots, \metav{A}_n$ 证明出一个矛盾\\
\\
$\metav{A}_1, \metav{A}_2, \ldots, \metav{A}_n \therefore \metav{C}$ 是\textbf{有效的}吗？
& 给出一个以 $\metav{A}_1, \metav{A}_2, \ldots, \metav{A}_n$ 为假设、以 \metav{C} 为结论的证明
& 给出一个使得每个 $\metav{A}_1, \metav{A}_2, \ldots, \metav{A}_n$ 为真而 \metav{C} 为假的解释\\
\end{tabular*}
\end{center}
\ifHTMLtarget
\end{table}


\else
\end{sidewaystable}
\fi